\documentclass[12pt]{article}
\usepackage[a4paper,margin=1in]{geometry}\usepackage[T1]{fontenc}
\usepackage{lmodern,microtype}\usepackage{amsmath,amssymb,amsthm,mathtools}
\usepackage{booktabs,array,enumitem,xcolor}\usepackage[numbers,sort&compress]{natbib}
\usepackage[colorlinks=true,linkcolor=blue!55!black,citecolor=blue!55!black,urlcolor=blue!55!black]{hyperref}
\linespread{1.07}
\newtheorem{theorem}{Theorem}[section]\newtheorem{proposition}[theorem]{Proposition}
\newtheorem{corollary}[theorem]{Corollary}
\theoremstyle{remark}\newtheorem{remark}[theorem]{Remark}
\newcommand{\norm}[1]{\left\lVert#1\right\rVert}

\title{Transverse Quenching in Quartet Coefficient Space\\
\large Exact Anisotropy Near the Degenerate Axial Boundary}
\author{Liang Wang}\date{July 2026}
\begin{document}\maketitle
\begin{abstract}
The centered conjugate-quartet manifold has axial coordinate $u$ and
asymmetry coordinate $\eta$.  RH-216 proves that $u\to1$ forces the
polynomial toward $(z^2-1)^2$ uniformly in $\eta$.  We sharpen that statement
by computing the full coefficient Jacobian and proving an exact transverse
Lipschitz bound.

For
\[
 \Phi(u,\eta)=(2-4u,\ 4\sqrt u(1-u)\eta,\
                 1-\eta^2(1-u)^2),
\]
the $\eta$ derivative satisfies
\[
 \sup_{|\eta|\le1}\norm{\partial_\eta\Phi(u,\eta)}_2
 \le 2(1-u)\sqrt{4u+(1-u)^2}.
\]
The same constant controls every finite transverse chord.  It vanishes
linearly as $u\to1$.  In contrast,
$\norm{\partial_u\Phi}\ge4$ because the first component is always $-4$.
Thus coefficient space becomes intrinsically anisotropic: axial motion
remains visible while transverse motion is quenched.

On the sixteen-level physical atlas, the maximum channelwise
transverse-to-axial Jacobian ratio decreases from $0.3752$ at
$\sigma=0.02$ to $0.2389$ at $\sigma=0.00125$.  The finest left/right
transverse coefficient leg is only $0.00136$.  These finite values support
the geometric mechanism but do not prove $u\to1$, convergence of $\eta$, or
a one-dimensional physical flow.
\end{abstract}

\section{Why a narrow corridor may become coefficientwise invisible}

RH-214 observes that $\eta$ fluctuates in a mature interval of width below
$0.038$, while $u$ increases across every sampled transition
\citep{WangRH214}.  RH-216 shows that the entire $u=1$ edge of the shape
rectangle maps to one polynomial \citep{WangRH216}.  These facts suggest a
specific mechanism: the coefficient map may contract transverse differences
as the axial coordinate grows.

This paper proves that mechanism directly.  It is a statement about the
exact quotient geometry of RH-213, independent of any statistical fit or
physical limit \citep{WangRH213}.

\section{Coefficient map}

Discarding the fixed leading and cubic coefficients, define
\begin{equation}\label{eq:phi}
 \Phi(u,\eta)=
 \begin{pmatrix}
 c_2\\c_3\\c_4
 \end{pmatrix}
 =\begin{pmatrix}
 2-4u\\
 4\sqrt u(1-u)\eta\\
 1-\eta^2(1-u)^2
 \end{pmatrix},
\end{equation}
for $0<u\le1$ and $|\eta|\le1$.

\begin{proposition}[Exact Jacobian]\label{prop:jacobian}
For $u>0$,
\begin{equation}\label{eq:jacobian}
 D\Phi(u,\eta)=
 \begin{pmatrix}
 -4 & 0\\[2mm]
 \displaystyle\frac{2\eta(1-3u)}{\sqrt u}
     &4\sqrt u(1-u)\\[3mm]
 2\eta^2(1-u)&-2\eta(1-u)^2
 \end{pmatrix}.
\end{equation}
\end{proposition}
\begin{proof}
Differentiate the three components of \eqref{eq:phi}.  In particular,
\begin{equation}
 \frac{d}{du}\bigl[\sqrt u(1-u)\bigr]
 =\frac{1-3u}{2\sqrt u}.
\end{equation}
\end{proof}

The singularity in the displayed $u$ derivative at $u=0$ reflects loss of
the positive/negative real branch label there.  The physical mature region
stays away from that boundary.

\section{Uniform transverse bound}

\begin{theorem}[Transverse Lipschitz quenching]\label{thm:quench}
For fixed $u\in[0,1]$ and any $\eta_1,\eta_2\in[-1,1]$,
\begin{equation}\label{eq:lipschitz}
 \norm{\Phi(u,\eta_1)-\Phi(u,\eta_2)}_2
 \le L_\eta(u)|\eta_1-\eta_2|,
\end{equation}
where
\begin{equation}\label{eq:L}
 \boxed{L_\eta(u)=2(1-u)\sqrt{4u+(1-u)^2}.}
\end{equation}
Consequently $L_\eta(u)=O(1-u)$ as $u\to1$.
\end{theorem}
\begin{proof}
At fixed $u$, the coefficient difference is
\begin{equation}
 \begin{pmatrix}
 0\\
 4\sqrt u(1-u)(\eta_1-\eta_2)\\
 -(1-u)^2(\eta_1+\eta_2)(\eta_1-\eta_2)
 \end{pmatrix}.
\end{equation}
Since $|\eta_1+\eta_2|\le2$, its squared norm is at most
\begin{equation}
 4(1-u)^2\bigl[4u+(1-u)^2\bigr]|\eta_1-\eta_2|^2.
\end{equation}
Taking square roots proves the estimate.  The asymptotic order follows
because the square-root factor tends to two.
\end{proof}

This is a global chord estimate, not merely a local derivative statement.
It remains valid at $u=0$ and on the $\eta$ edges.

\section{Axial sensitivity does not vanish}

\begin{proposition}[Persistent axial visibility]\label{prop:axial}
For every $u>0$ and $|\eta|\le1$,
\begin{equation}\label{eq:axiallower}
 \norm{\partial_u\Phi(u,\eta)}_2\ge4.
\end{equation}
\end{proposition}
\begin{proof}
The first component of $\partial_u\Phi$ is $-4$.
\end{proof}

Combining Theorem~\ref{thm:quench} and
Proposition~\ref{prop:axial} gives the uniform anisotropy ratio
\begin{equation}\label{eq:ratio}
 \frac{\sup_{\eta}\norm{\partial_\eta\Phi(u,\eta)}_2}
      {\inf_{\eta}\norm{\partial_u\Phi(u,\eta)}_2}
 \le \frac{1-u}{2}\sqrt{4u+(1-u)^2}\longrightarrow0.
\end{equation}

Thus the map itself, not a fitted trajectory, singles out the axial
direction near $u=1$.

\section{Coefficient quenching versus root quenching}

The imaginary root heights are
\begin{equation}
 b=\sqrt{(1-u)(1+\eta)},\qquad
 d=\sqrt{(1-u)(1-\eta)}.
\end{equation}
Their $\eta$ derivatives can become large near $|\eta|=1$ because of the
square root.  Therefore Theorem~\ref{thm:quench} is specifically a
coefficient-space theorem.  Symmetric polynomials cancel and square the
transverse root motion, producing the stronger $O(1-u)$ coefficient scale.

This distinction matters for the divisor-first route.  A scalar determinant
can stabilize even when individual branch coordinates are less regular, just
as RH-210 showed that divisor stability need not imply projector stability
\citep{WangRH210}.

\section{Exact path decomposition between channels}

At one scale let the left and right shapes be $(u_L,\eta_L)$ and
$(u_R,\eta_R)$.  Insert the corner $(u_R,\eta_L)$:
\begin{equation}\label{eq:path}
 \Phi(u_R,\eta_R)-\Phi(u_L,\eta_L)
 =\underbrace{\Phi(u_R,\eta_L)-\Phi(u_L,\eta_L)}_{\text{axial leg}}
 +\underbrace{\Phi(u_R,\eta_R)-\Phi(u_R,\eta_L)}_{\text{transverse leg}}.
\end{equation}
This is an exact telescoping identity.  The transverse leg obeys
\begin{equation}
 \norm{\text{transverse leg}}_2
 \le L_\eta(u_R)|\eta_R-\eta_L|.
\end{equation}

The frozen audit evaluates both legs at all sixteen scales.  Eight hundred
random chord tests have no positive bound violation, and all path residuals
are machine zero.

\section{Finite physical anisotropy}

The exact Jacobian is evaluated at each of the 32 physical shape points.
Representative mature values are summarized by
\begin{center}
\begin{tabular}{lrr}
\toprule
diagnostic & $\sigma=0.02$ & $\sigma=0.00125$\\
\midrule
maximum transverse/axial Jacobian ratio & $0.37516$ & $0.23883$\\
maximum channel $|\Delta\eta|$ & $0.00145$ & $0.00141$\\
maximum transverse channel-leg norm & small & $0.001354$\\
\bottomrule
\end{tabular}
\end{center}

The ratio decreases as the observed axial clock advances, in agreement with
\eqref{eq:ratio}.  It is not yet extremely small: $u\approx0.718$ remains far
from the degenerate edge.  The numbers therefore support the mechanism
without demonstrating an asymptotic regime.

\section{Implications and limitations}

If a future theorem proves $u_\sigma\to1$, then any bounded $\eta_\sigma$
sequence---convergent or not---has vanishing transverse coefficient effect.
The coefficient path could become asymptotically one-dimensional even while
the transverse coordinate continues to oscillate.

The converse is not supplied here.  Finite decline of the sensitivity ratio
does not prove $u\to1$, and coefficient anisotropy does not identify a
scale-independent map advancing $u$.  RH-218 tests that stronger recurrence
claim separately.

No growing root cloud, locally uniform determinant, Gate-A closure,
Hilbert--P\'olya operator, arithmetic trace formula, or zeta identification
is obtained.

\bibliographystyle{plainnat}\bibliography{references}\end{document}
